% se puede agregar la opción [english] para 
%  memorias o tesis en inglés
\documentclass{umemoria} 

\usepackage[utf8]{inputenc}

\depto{Departamento de ???}
\author{Nombre Completo Autor}
\title{Título de la Memoria/Tesis}

% incluir ambos comandos para una doble titulación
%  o quitar el comando que no aplica
\memoria{Ingenier[o/a?] Civil en ???}
\tesis{Magister en ???}
% \tesis{Doctor en ???} % incluir solo este comando para doctorados

% puede haber varios profesores guía seperados por coma;
% pero se es una memoria, solo puede haber un profesor guía
\guia{Nombre Completo Guía} 

% puede haber varios profesores co-guía seperados por coma;
% pero si es una memoria, el profesor co-guía será el primer
% integrante de la comisión
% \coguia{Nombre Guía} % incluir en caso de co-guía de *tesis*

%\cotutela{Nombre Institución} % incluir en caso de cotutela

\comision{Nombre Completo Uno,Nombre Completo Dos,Nombre Completo Tres}
%\auspicio{Nombre Institución} % incluir en caso de recibir financiamiento

% tiene que ser el año en que se da el examen de grado (defensa)
%\anho{2021} % incluir solo para reemplazar el año actual

\usepackage{lipsum}

\begin{document}

\frontmatter
\maketitle

\begin{resumen}
\lipsum[1-4]
\end{resumen}

% opcional: incluir para tesis en inglés;
%  en general hay que tener el resumen y abstract
%   en ambos idiomas
%\begin{abstract}
%\lipsum[1-4]
%\end{abstract}

\begin{dedicatoria}
Una dedicatoria corta.
\end{dedicatoria}

\begin{thanks}
\lipsum[1-2]
\end{thanks}

\tableofcontents
\listoftables % opcional
\listoffigures % opcional

\mainmatter

\begin{intro}
\lipsum[30-35]
\begin{enumerate}
	\item Item 1
	\begin{enumerate}
		\item Subitem 1
		\item Subitem 2 (ver Figura \ref{logofcfm})
	\end{enumerate}
	\item Item 2
	\item Item 3
\end{enumerate}
\begin{teo}
Se tiene que $$\int_0^t e^sds=e^t-1.$$
\end{teo}
\lipsum[36-40]
\end{intro}
\chapter{Segundo}

\lipsum[1-3]

\begin{defn}[ver \cite{KAR00}] Definición definitiva $$\frac{d}{dx}\int_a^xf(y)dy=f(x).$$\end{defn}
\chapter{Tercero}
\lipsum[50-60]
\chapter{Conclusión}

\lipsum[130-132]
\begin{figure}[tbh]
	\centering
	\includegraphics[scale=.2]{imagenes/fcfm.pdf}
	\caption{Logo de la Facultad}
	\label{logofcfm}
\end{figure}

\lipsum[133-134]

\begin{table}[tbh]
	\centering
	\caption{Tabla 1}
	\label{tabla:1}
	\begin{tabular}{|c|c|r|}
		\hline
		\textbf{Campo 1} & \textbf{Campo 2} & \textbf{Num} \\\hline
		Valor 1a & Valor 2a & 3\\
		Valor 1b & Valor 2b & 3\\
		\hline
	\end{tabular}

\end{table}

\lipsum[135]


% ver https://www.overleaf.com/learn/latex/Glossaries
% \input{glosario.tex} % opcional

\nocite{*}
\bibliographystyle{plain}
\bibliography{bibliografia}

% opcional ...
%\appendix
%\chapter{Anexo}
\lipsum[50-60]
\end{document}
